\documentclass[11pt]{article}
\usepackage[a4paper,margin=1in]{geometry}
\usepackage[T1]{fontenc}
\usepackage[utf8]{inputenc}
\usepackage{lmodern}
\usepackage{hyperref}
\usepackage{longtable}
\usepackage{enumitem}
\usepackage{titlesec}

\title{DNA String Matcher\\Detailed Technical Stack and Architecture Report}
\author{Project Documentation Draft}
\date{April 14, 2026}

\begin{document}
\maketitle
\tableofcontents
\newpage

\section{Project Overview}
\textbf{Project Name:} DNA Pattern Matcher / GeneFlow DNA String Matcher.

\textbf{Goal:} Detect DNA subsequences inside larger genome strings using a Deterministic Finite Automaton (DFA) based matcher, with both desktop and web user interfaces.

\textbf{Academic Framing:}
\begin{itemize}[leftmargin=*]
  \item Practical application of automata theory to bioinformatics pattern search
  \item DFA construction, transition traversal, and complexity discussion
  \item Multi-interface software implementation (desktop + web)
\end{itemize}

\section{Technology Stack}
\subsection{Language and Core Runtime}
\begin{itemize}[leftmargin=*]
  \item Python
  \item Documented version references vary across files (3.9+, 3.10+, and 3.13 noted in docs/scripts)
\end{itemize}

\subsection{Algorithm Layer}
\begin{itemize}[leftmargin=*]
  \item Custom DFA transition table builder
  \item Custom DFA-based matcher and transition tracer
\end{itemize}

\subsection{Desktop Stack}
\begin{itemize}[leftmargin=*]
  \item PyQt6
  \item PyQt6-WebEngine
  \item Matplotlib embedded in Qt
\end{itemize}

\subsection{Web Stack}
\begin{itemize}[leftmargin=*]
  \item Streamlit
  \item Plotly
  \item Pandas
\end{itemize}

\subsection{Data and Bioinformatics}
\begin{itemize}[leftmargin=*]
  \item BioPython (FASTA parsing via SeqIO)
  \item Custom genome loader with DNA alphabet filtering
\end{itemize}

\subsection{AI Integration (Optional)}
\begin{itemize}[leftmargin=*]
  \item Hugging Face Router endpoint for chat completions
  \item Multi-model fallback strategy
  \item Environment and .env token resolution support
\end{itemize}

\section{Declared Dependencies}
\begin{longtable}{|p{0.45\textwidth}|p{0.45\textwidth}|}
\hline
\textbf{Dependency} & \textbf{Role} \\
\hline
PyQt6>=6.7 & Desktop UI \\
\hline
PyQt6-WebEngine>=6.7 & Embedded web rendering in desktop flow \\
\hline
matplotlib>=3.7 & Charts and visual exports \\
\hline
google-generativeai>=0.3.0 & Legacy/alternate AI dependency reference \\
\hline
biopython>=1.83 & FASTA parsing \\
\hline
streamlit>=1.28 & Web UI \\
\hline
plotly>=5.17 & Interactive web charts \\
\hline
pandas>=2.0 & Tabular processing for web app \\
\hline
certifi>=2024.2.2 & SSL certificates for API networking \\
\hline
\end{longtable}

\textbf{Consistency note:} Active AI code uses Hugging Face HTTP endpoints, while requirements and some docs still mention Google/Gemini related setup.

\section{Architecture}
\subsection{Layered View}
\begin{enumerate}[leftmargin=*]
  \item \textbf{Input Layer:} Pattern text, genome text, FASTA file, optional natural-language query.
  \item \textbf{Processing Layer:} DFA construction, genome scan, accept-state detection, optional transition trace.
  \item \textbf{Presentation Layer:} PyQt desktop visualization and Streamlit web visualization.
  \item \textbf{Support Layer:} AI helper module, test scripts, verification script.
\end{enumerate}

\subsection{Entry Points}
\begin{itemize}[leftmargin=*]
  \item main.py: launcher menu
  \item app\_desktop.py: desktop app start
  \item app\_web.py: Streamlit web app
  \item app\_streamlit.py: alternate Streamlit app
  \item demo.py: CLI demo script (contains stale references)
\end{itemize}

\section{Repository Structure Notes}
\begin{itemize}[leftmargin=*]
  \item Root contains primary implementation modules.
  \item pyqt-app/ contains a mirrored copy of much of the project.
  \item assets/ stores UI styles.
  \item outputs/ stores generated visual files.
  \item web/ is present but empty.
\end{itemize}

\section{Algorithm Details}
\subsection{DFA Construction}
Given pattern length $m$, states are represented as $0..m$. For each state and each alphabet symbol in \{A,T,C,G\}, the transition function computes the largest prefix length matched after consuming the symbol.

\subsection{Matching Procedure}
Given genome length $n$, the matcher performs a single left-to-right scan. State updates follow the transition table. Whenever accept state $m$ is reached, a match start index is recorded.

\subsection{Complexity Discussion}
\begin{itemize}[leftmargin=*]
  \item Scan phase: $O(n)$
  \item Preprocessing depends on pattern and transition construction method
  \item Project framing emphasizes linear-time DFA traversal behavior
\end{itemize}

\section{Data Validation and Handling}
\begin{itemize}[leftmargin=*]
  \item Pattern validation enforces alphabet \{A,T,C,G\}.
  \item Genome input is uppercased and sanitized to valid bases.
  \item FASTA parsing supported in web flow via BioPython and in desktop via local loader.
  \item Primary reported matches are 1-indexed positions.
\end{itemize}

\section{UI and Visualization}
\subsection{Desktop UI}
\begin{itemize}[leftmargin=*]
  \item Custom-painted DFA state graph and transitions
  \item Timeline scanner with current index and match highlighting
  \item Optional stepwise animation from transition traces
  \item Embedded Matplotlib plotting
\end{itemize}

\subsection{Web UI}
\begin{itemize}[leftmargin=*]
  \item Streamlit interface with custom HTML/CSS sections
  \item Plotly charts and metric displays
  \item Session-state based analysis workflow
\end{itemize}

\section{AI Subsystem}
\begin{itemize}[leftmargin=*]
  \item Natural-language to DNA-pattern extraction
  \item Friendly error handling for auth/quota/model/network issues
  \item Model fallback list for reliability
  \item Optional motif suggestion logic via k-mer frequency counting
\end{itemize}

\section{Testing and Verification}
\subsection{test\_suite.py}
Covers functional cases, overlapping matches, biological motif examples, edge conditions, and performance checks.

\subsection{verify\_system.py}
Covers runtime dependency verification, algorithm smoke tests, and project file presence checks.

\subsection{Current QA Style}
Script-based testing is present; no formal pytest-based automated CI pipeline is detected.

\section{Known Inconsistencies and Technical Debt}
\begin{enumerate}[leftmargin=*]
  \item AI dependency/documentation mismatch (Google/Gemini mentions vs Hugging Face implementation).
  \item demo.py contains stale imports/references (aho\_corasick, visualize, app.py).
  \item Duplicated codebase in pyqt-app/ increases maintenance overhead.
  \item Python minimum version statements are inconsistent across files.
  \item Desktop plotting backend choice may need verification with PyQt6 runtime settings.
\end{enumerate}

\section{Deployment and Execution}
\subsection{Typical Local Run}
\begin{enumerate}[leftmargin=*]
  \item Create/activate virtual environment
  \item Install requirements
  \item Run main.py and choose desktop or web mode
\end{enumerate}

\subsection{Direct Commands}
\begin{itemize}[leftmargin=*]
  \item Web: streamlit run app\_web.py
  \item Desktop: python3 app\_desktop.py
\end{itemize}

\section{Presentation-Ready Structure (For PPT)}
\begin{enumerate}[leftmargin=*]
  \item Problem and Motivation
  \item Automata Theory Mapping
  \item Full Architecture
  \item Tech Stack by Layer
  \item DFA Logic and Complexity
  \item Interface Demo (Desktop + Web)
  \item AI-Assist Workflow
  \item Testing Results
  \item Gaps and Improvements
  \item Roadmap
\end{enumerate}

\section{Recommended Roadmap}
\begin{enumerate}[leftmargin=*]
  \item Consolidate root and pyqt-app duplicates
  \item Standardize AI provider and cleanup dependency list
  \item Repair or deprecate stale demo paths
  \item Introduce pytest and CI automation
  \item Unify Python version policy
\end{enumerate}

\section{Conclusion}
The project demonstrates a strong automata-driven DNA matching foundation with rich visualization and optional AI assistance. It is presentation-ready for academic reporting, with clear opportunities to improve maintainability and stack consistency before production-level release.

\end{document}
